% Generated draft: requires mathematical and source review.
\documentclass[11pt]{article}
\usepackage[margin=25mm]{geometry}
\usepackage{fontspec}
\setmainfont{Latin Modern Roman}
\usepackage{amsmath,amssymb,mathtools,unicode-math}
\setmathfont{Latin Modern Math}
\usepackage{enumitem,needspace,xurl,hyperref}
\setlength{\parindent}{0pt}
\setlength{\parskip}{5pt}
\setlength{\emergencystretch}{3em}
\newcommand{\N}{\mathbb{N}}
\newcommand{\Z}{\mathbb{Z}}
\newcommand{\Q}{\mathbb{Q}}
\newcommand{\R}{\mathbb{R}}
\newcommand{\C}{\mathbb{C}}
\newcommand{\F}{\mathbb{F}}
\newcommand{\A}{\mathbb{A}}
\newcommand{\PP}{\mathbb{P}}
\newcommand{\E}{\mathbb{E}}
\newcommand{\eps}{\varepsilon}
\newcommand{\dd}{\,\mathrm{d}}
\begin{document}
\section*{UnsolvedMath ID: 30005804}
\textbf{Source problem code:} catalogue entry
\subsection{Definitions and mathematical statement}
For a compact convex set \(K\subseteq\R^2\) and a vector \(t\in\R^2\), its translate is \(K+t=\{x+t:x\in K\}\). A finite family \(\mathcal F\) is pierced by at most three points if there exists a set \(P\subseteq\R^2\) with \(|P|\leq 3\) such that \(P\cap A\neq\varnothing\) for every \(A\in\mathcal F\). The piercing points need not belong to \(K\). An empty family is pierced vacuously.

\textbf{Dol'nikov's colorful transversal conjecture.} For every compact convex set \(K\subseteq\R^2\) and every three finite families \(\mathcal F_1,\mathcal F_2,\mathcal F_3\) whose members are translates of this same set \(K\), assume that
\[
\forall i,j\in\{1,2,3\},\quad i\neq j\ \Longrightarrow\ \bigl(\forall A\in\mathcal F_i\ \forall B\in\mathcal F_j,\ A\cap B\neq\varnothing\bigr).
\]
Then
\[
\exists j\in\{1,2,3\}\ \exists P\subseteq\R^2,\qquad |P|\leq 3\quad\text{and}\quad \forall A\in\mathcal F_j,\ P\cap A\neq\varnothing.
\]
No intersection condition is imposed on two members of the same family. No nonempty-interior assumption on \(K\), or nonemptiness assumption on the families, is added. Here ``three points'' is expressed as ``at most three points''; any smaller piercing set can be enlarged to three distinct points in \(\R^2\).

\textbf{Source qualifications.} This is the conjecture stated in the supplied clean record. The supplied original excerpt is truncated and contains apparently interleaved, corrupted material concerning projection theory and the colorful Helly theorem; those passages do not supply additional hypotheses for this statement. The record labels the problem partially solved and reports special cases for centrally symmetric bodies, triangles, disks, bodies of constant width, and bodies sufficiently close to a disk. These are supplied literature assessments, not independently verified mathematical conclusions here. Neither the meaning of ``sufficiently close'' nor the hypotheses of the cited special-case results are specified in the record.
\subsection{Short English statement}
Dol'nikov's conjecture asks whether three finite families of translates of one compact convex planar set, with every pair from different families intersecting, must have one family that can be met by a common set of at most three points. The supplied record reports partial results; this statement-only draft requires human mathematical and source review.
\subsection{Sources}
\begin{itemize}
\item Local registry: \texttt{data/unsolvedmath/problems.json}, permanent UnsolvedMath ID \(30005804\), title ``Dolnikov's Colorful Transversal Conjecture,'' exact source problem code \texttt{catalogue entry}. The supplied record labels its statement \texttt{corrected\_verified} and says it was reconstructed and checked against the official report; this verification claim is supplied metadata and was not independently checked here.
\item \emph{Discrete Geometry} (2024), Oberwolfach Reports, pp.\ 137--202, DOI \texttt{10.4171/owr/2024/3}: \url{https://doi.org/10.4171/owr/2024/3}. This is the supplied authoritative report citation. The excerpt attributes the conjecture to Dol'nikov and identifies the presentation ``On Dol'nikov's Conjecture'' with Edgardo Rold\'an-Pensado, joint work with Cuauht\'emoc Gomez-Navarro and Leonardo Mart\'inez-Sandoval; it does not identify these speakers as the conjecture's originators.
\item \emph{Transversals to Colorful Intersecting Convex Sets}: \url{https://doi.org/10.1007/s00454-024-00669-3}. Supplied literature-source link; no author list or detailed theorem statement is supplied for this entry.
\item \emph{On a colorful problem by Dol'nikov concerning translates of convex bodies}: \url{https://arxiv.org/abs/2307.07714}. Supplied literature-source link; no detailed theorem statement is supplied for this entry.
\item Historical attribution in the supplied excerpt: Jer\'onimo-Castro, Magazinov and Sober\'on are credited with the centrally symmetric and triangle cases in 2015 and a stronger result for a ``circle.'' Further progress in 2023 is attributed to collaboration with Gomez-Navarro, but the excerpt ends before specifying the additional sets. Its numbered references [1] and [2] are not accompanied by their bibliography entries.
\item Literature assessment checked at \texttt{2026-08-22}: the supplied curation labels the problem partially solved, reports the special cases listed above, and says that the full compact-convex-set statement remains open. This is supplied dated curation derived from the OWR material, not an independently verified assessment; the record explicitly requests independent verification of time-sensitive status and cited claims.
\end{itemize}
\end{document}
